% Preference-conditioned pick-and-place: VLA baseline vs. full VLM--VLA system.
% Fragment for \input{} into the thesis. Needs: booktabs, tabularx (both standard).
% writeup_standalone.tex wraps it for compiling on its own.

\section{Preference-conditioned pick-and-place: VLA baseline vs.\ full VLM--VLA system}
\label{sec:pref-eval}

\subsection{Aim}

A VLA policy fine-tuned on single-object instructions (``put the banana on the tray'') is
given instructions that carry a \emph{preference}: an exclusion, an order, a condition on
the objects, or a count. This experiment measures whether the VLA can follow such
preferences on its own. It then measures how much is gained by placing it inside the full
system, where a VLM decomposes the preference into the base sub-tasks the VLA was trained
on, monitors each grasp, evaluates every sub-task and the final scene, and retries.

The experiment asks four questions:
\begin{enumerate}
  \item \textbf{Intent:} does the system go for the objects the preference allows, in the
    order it requires, and stop where it should?
  \item \textbf{Execution:} does the final tray match the preference?
  \item \textbf{Failure sources:} which component causes the remaining failures?
  \item \textbf{Ablation:} does in-round replanning after a failed sub-task help?
\end{enumerate}

\subsection{Method}

\subsubsection{Setup}
\begin{itemize}
  \item \textbf{Robot:} SO-101 arm, a wrist camera, and a top-down RealSense camera. The
    VLA is \texttt{ant0nh/pi05\_500\_30k} in both conditions.
  \item \textbf{Scene:} a pink tray and three objects: a banana, a brown plush moose toy
    and a brown cylindrical pouch.
  \item \textbf{Layouts:} objects sit in three slots, ordered nearest-first (front, back,
    left). Three layouts (\texttt{ptb}, \texttt{tbp}, \texttt{bpt}) make each object the
    nearest once, so a pure ``take the nearest object'' policy can be separated from
    following the instruction.
\end{itemize}

\subsubsection{Prompt set and worklist}
The baseline's prompt set is summarised in Table~\ref{tab:pref-prompts}. It gives
\textbf{57 trials}. Each prompt cycles through the layouts, and the ``one object'' and
``two objects'' prompts get five runs each. The queue is shuffled with seed~0. The full
system ran the \emph{same} 57 (prompt, layout) items as the baseline, so every system
trial is paired with a baseline trial.

\begin{table}[t]
  \centering
  \caption{Prompt set and number of trials per category.}
  \label{tab:pref-prompts}
  \begin{tabularx}{\linewidth}{lrrX}
    \toprule
    Category & Prompts & Trials & Example \\
    \midrule
    Omission        & 6 & 12 & ``put everything on the tray except the banana and the toy'' \\
    Sequencing      & 5 & 10 & ``put the toy on the tray first, then the pouch, then the banana'' \\
    Conditional     & 6 & 12 & ``put everything that is not edible on the tray'' \\
    Quantity        & 4 & 14 & ``put two objects on the tray'' \\
    Prior (control) & 1 &  6 & ``put everything on the tray'' \\
    Base (control)  & 3 &  3 & ``put the pouch on the tray'' \\
    \midrule
    Total           & 25 & 57 & \\
    \bottomrule
  \end{tabularx}
\end{table}

\subsubsection{Conditions}
\paragraph{Baseline.} The raw preference prompt is sent straight to the VLA, with no VLM
and no monitor, for one 60\,s episode.

\paragraph{Full system.} Each prompt runs through the orchestrator. The VLM is
Qwen3-VL-8B at temperature 0.1.
\begin{enumerate}
  \item \textbf{Decompose:} a two-pass planner identifies the visible objects, filters out
    fixtures, then decomposes the prompt into ``put the X on the tray'' sub-tasks.
  \item \textbf{Execute:} one VLA episode per sub-task, up to 100\,s.
  \item \textbf{Live monitor:} at the grasp, the arm is frozen; a two-pass reach check
    runs, and the wrist camera verifies the held object. A confirmed stop ends the
    episode; the same sub-task is retried once, and a second stop skips it.
  \item \textbf{Sub-task evaluation:} a 32-frame clip of the episode is judged by the VLM.
    If the arm did not move, a before/after pre-check runs first. One retry is allowed.
  \item \textbf{Final evaluation:} a before/after scene survey, the instruction's
    requirements (objects on or off the tray, counts), and a comparison in code. If the
    verdict is incomplete, one retry round re-plans over the objects not yet on the tray,
    with a count cap for count prompts.
\end{enumerate}

\paragraph{Replanning ablation.} In-round replanning after a sub-task fails all its
attempts was enabled for the first pass through the worklist. The replan fired in 11
trials. Those 11 (prompt, layout) items were re-run with replanning disabled, which is the
final configuration.
\begin{itemize}
  \item \textbf{Replan ON} is the first pass.
  \item \textbf{Replan OFF} is the 46 first-pass trials that never replanned (equivalent
    to replanning off) plus the 11 re-runs.
\end{itemize}

\subsubsection{Labels and metrics}
After every trial the operator recorded two sequences over the whole task, all episodes
included:
\begin{itemize}
  \item \textbf{attempted:} the objects the arm went for, in order (a real approach and
    grasp attempt, held or not);
  \item \textbf{placed:} the objects on the tray at the end, in order.
\end{itemize}
Both conditions are scored with the same scoring function; the VLM's own verdicts are
never the score.
\begin{itemize}
  \item \textbf{Intent} (from \emph{attempted}): \emph{reach admissible}, the first attempt
    is allowed by the preference; \emph{intent complete}, the attempted sequence satisfies
    the whole preference; \emph{intent violated}, it attempts an excluded or unselected
    object, breaks the order, or exceeds the count.
  \item \textbf{Execution} (from \emph{placed}): \emph{first placement admissible},
    \emph{complete} (the placed sequence satisfies the preference), and \emph{violated}.
  \item \textbf{Prior lift:} reach admissible minus what the prior (``put everything'')
    trials' first-attempt distribution predicts for the same allowed set, by proximity
    rank.
\end{itemize}
Paired comparisons use McNemar's exact test; intervals are 95\% Wilson intervals.

\subsection{Results}

\subsubsection{Headline}
Table~\ref{tab:pref-headline} gives the results over all 57 paired trials, controls
included. On the 48 preference trials alone (the 9 prior and base controls excluded), task
complete is 10/48 (21\%), 31/48 (65\%) and 38/48 (79\%) for the baseline, replan ON and
replan OFF, and violated is 26/48 (54\%), 12/48 (25\%) and 7/48 (15\%). These are the
sums of the four preference rows of Table~\ref{tab:pref-category}; the headline table adds
the controls (3, 8 and 8 complete).

\begin{table}[t]
  \centering
  \caption{Results over all 57 paired trials, controls included. Brackets give 95\%
    Wilson intervals.}
  \label{tab:pref-headline}
  \begin{tabular}{lccc}
    \toprule
    Metric & Baseline & System, replan ON & System, replan OFF \\
    \midrule
    First attempt admissible   & 33 (58\%) & 55 (96\%) & 55 (96\%) \\
    Intent complete            & 10 (18\%) & 43 (75\%) & \textbf{49 (86\%)} \\
    Intent violated            & 34 (60\%) & 13 (23\%) & \textbf{8 (14\%)} \\
    First placement admissible & 27 (47\%) & 53 (93\%) & 54 (95\%) \\
    \midrule
    Task complete & 13 (23\%) [14--35] & 39 (68\%) [56--79] & \textbf{46 (81\%) [69--89]} \\
    Violated      & 26 (46\%) [33--58] & 13 (23\%) [14--35] & \textbf{8 (14\%) [7--25]} \\
    \bottomrule
  \end{tabular}
\end{table}

Paired against the baseline (replan OFF):
\begin{itemize}
  \item \textbf{Complete:} 35 trials only the system completed, 2 only the baseline
    completed ($p \approx 10^{-8}$).
  \item \textbf{Violated:} 22 trials only the baseline violated, 4 only the system
    violated ($p \approx 5 \times 10^{-4}$).
  \item \textbf{Intent complete:} 40 vs.\ 1 ($p \approx 4 \times 10^{-11}$).
\end{itemize}

\paragraph{Operator aborts.} Three replan-OFF trials were aborted by the operator. In both
``put the yellow object on the tray'' trials (\#54, \#58) the planner selected all three
objects; the operator aborted after the banana was placed, which left a ``complete''
tray. Scored on what the system planned, these two are failures. Counted that way, replan
OFF completes 44/57 (77\%) overall and 36/48 (75\%) of the preference trials.

\subsubsection{By category}
Table~\ref{tab:pref-category} breaks the results down by category.

\begin{table}[t]
  \centering
  \caption{Results by category. Complete and violated are given as
    baseline\,/\,replan ON\,/\,replan OFF; intent complete as baseline\,/\,replan OFF.}
  \label{tab:pref-category}
  \begin{tabular}{lrccc}
    \toprule
    Category & $n$ & Complete & Violated & Intent complete \\
    \midrule
    Omission    & 12 & 0\,/\,10\,/\,11 & 12\,/\,2\,/\,1 & 0\,/\,11 \\
    Sequencing  & 10 & 0\,/\,3\,/\,5   & 6\,/\,4\,/\,3  & 0\,/\,8 \\
    Conditional & 12 & 2\,/\,7\,/\,10  & 6\,/\,5\,/\,2  & 2\,/\,9 \\
    Quantity    & 14 & 8\,/\,11\,/\,12 & 2\,/\,1\,/\,1  & 5\,/\,13 \\
    Prior       &  6 & 0\,/\,6\,/\,6   & 0\,/\,0\,/\,0  & 0\,/\,6 \\
    Base        &  3 & 3\,/\,2\,/\,2   & 0\,/\,1\,/\,1  & 3\,/\,2 \\
    \bottomrule
  \end{tabular}
\end{table}

\subsubsection{Did the VLA respond to the preference?}
\begin{itemize}
  \item \textbf{Baseline:} the first attempt was admissible in 24 of 48 preference
    trials (50\%; 58\% over all 57 trials including controls), against 53\% predicted from the prior's proximity-rank distribution: a lift of
    $-3$ points. Its first attempt was the nearest object in 65\% of trials. When the
    prompt named an object, the first attempt was that object in 79\% of trials (keyword
    latching), including the excluded object on omission prompts. All 12 baseline
    omission trials were violations.
  \item \textbf{Full system:} 96\% admissible first attempts, a $+32$-point lift over the
    same prior. Nearest-first dropped to 21\%, so the choice of object no longer follows
    proximity.
\end{itemize}

\subsubsection{Cost}
Table~\ref{tab:pref-cost} compares the cost of the conditions. Oscillation, as measured
by the baseline's motion metric, was flagged in 2 of 153 system episodes and in about 4\%
of baseline preference trials.

\begin{table}[t]
  \centering
  \caption{Cost per trial.}
  \label{tab:pref-cost}
  \begin{tabular}{lccc}
    \toprule
     & Baseline & Replan ON & Replan OFF \\
    \midrule
    Wall time per trial (mean)    & 71\,s & 185\,s & 155\,s (median 144\,s) \\
    VLA episodes per trial (mean) & 1     & 3.25   & 2.68 \\
    Trials with a retry round     & --    & 19     & 16 \\
    Monitor stops                 & --    & 18     & 13 \\
    \bottomrule
  \end{tabular}
\end{table}

\subsubsection{Replanning ablation}
With replanning on, the replan fired in 11 of 57 trials, each after a sub-task was judged
failed on both attempts:
\begin{itemize}
  \item 4 times on the pouch, which was already on the tray (sub-task evaluator false
    fails);
  \item 5 times on an object name the VLA does not respond to, or on an invented object
    (``cylinder'', ``case'', ``pen case'', ``black object'');
  \item 2 times after a genuine grasp failure.
\end{itemize}
It never clearly recovered a trial. Six of the 11 ended violated:
\begin{itemize}
  \item on conditional prompts it added the unselected banana (\#15, \#22, \#41);
  \item it kept an invented or unknown name, which the VLA or the monitor then grounded to
    the wrong object (\#29 ``black object'', \#46 ``case'');
  \item \#47 broke the order.
\end{itemize}
It also planned three objects for ``one object'' (\#36), which ended incomplete.

With replanning off, 9 of the 11 re-runs completed. Only 3 of the 11 re-runs reached a
point where a replan would have fired. Of those, one changed (\#47, violated
$\rightarrow$ complete); \#44 completed and \#15 was violated in both conditions. The
overall ON $\rightarrow$ OFF gain (39 $\rightarrow$ 46 complete) is therefore mostly
variance from re-running the episodes, not a measured effect of replanning. The direct
evidence is that the replan caused six violations and recovered nothing.

\subsubsection{Component diagnostics (replan OFF)}
\begin{itemize}
  \item \textbf{Sub-task evaluator:} 35 of 137 judged attempts were failed. The
    no-movement pre-check rescued 21 attempts: the arm did not move because the object was
    already on the tray, nearly always the pouch after a false fail on the previous
    attempt (the pouch placed onto an occupied tray). The monitor's grasp freeze takes up
    most of the evaluation clip, and the placement appears only in the last few frames,
    brown on brown. Each pouch false fail costs a retry episode.
  \item \textbf{Final evaluator:} its verdict agreed with the labels in 41 of 54 trials
    (76\%). There were 10 false ``incomplete'' verdicts: a banana hidden under the stack
    read as ``missing'', or the pouch on a stack. Each triggers an unnecessary retry round.
    The 3 false ``complete'' verdicts are all sequencing prompts (\#12, \#32, \#52): the
    final evaluation ignores order by design.
  \item \textbf{Live monitor:} 13 stops in 6 trials. Correct wrong-object catches: \#13,
    \#35, \#49, and the banana stops in \#39. False stops: \#39's pouch, where it could not
    match ``pouch'' with ``pencil case'', and \#12, confused by the banana lying on the
    toy. \#45 was a wrong stop that happened to prevent a planned violation.
  \item \textbf{Retry rounds:} replaying the 19 saved round-2 frames offline, the live
    round-2 plans were correct in 9 of 19 and broke the prompt in 6. When told the task
    was ``judged incomplete'', the planner planned the leftover objects even when the
    prompt excluded them.
\end{itemize}

\subsubsection{Remaining failures}
Table~\ref{tab:pref-failures} lists the 11 replan-OFF trials that were not complete, 8 of
them violated, plus the two aborted ``yellow object'' trials, which are labelled complete
but count as failures under the stricter scoring. In \#15 the final evaluation's requirements decided the toy is not brown,
and in \#50 they did not recognise the ``pencil case'' as the excluded pouch.

\begin{table}[t]
  \centering
  \caption{Remaining failures with replanning off.}
  \label{tab:pref-failures}
  \begin{tabularx}{\linewidth}{Xlr}
    \toprule
    Source & Trials & Violated \\
    \midrule
    \multicolumn{3}{l}{\textit{Round-1 planner error}} \\
    \quad ``put the banana on the tray after the toy'' became one malformed sub-task
      (``put the banana after the moose toy on the tray''); only the banana was placed
      & \#32, \#52 & 2 \\
    \quad Two objects planned for ``one of the brown objects'' & \#20 & 1 \\
    \quad All three objects selected for ``the yellow object'' (operator aborted; labelled
      complete, not among the 11) & \#54, \#58 & (0) \\
    \midrule
    Final-evaluation misread; the retry round planned an excluded or unselected object
      & \#14, \#15, \#40, \#50 & 4 \\
    \midrule
    Grasp failure in a sequencing task; the object was placed later and the order broke
      & \#12 & 1 \\
    \midrule
    \multicolumn{3}{l}{\textit{Incomplete, nothing violated}} \\
    \quad Banana grasp failed; operator aborted & \#11 & 0 \\
    \quad Monitor repeatedly stopped the pouch grasp (``pouch'' vs.\ ``pencil case'')
      & \#39 & 0 \\
    \quad Round 2 planned nothing; the survey missed the toy & \#53 & 0 \\
    \bottomrule
  \end{tabularx}
\end{table}

\subsection{Discussion}

\paragraph{The system's gain comes from decomposition, not from better motor control.}
Both conditions use the same VLA. On the single-object base instructions it was trained
on, the VLA alone succeeds (3/3). On preference prompts it ignores the preference (lift
$-3$) and follows proximity or keyword latching. Omission makes this clearest: every
baseline omission trial placed an excluded object, while the system completed 11 of 12.
Decomposing the preference into the trained base instructions turns a
preference-following problem into one the VLA was trained to solve.

\paragraph{Where the system still fails, the cause is VLM perception, and the closed loop
amplifies it.} The sub-task evaluator's false fails, the final evaluator's false
``incomplete'' verdicts, and the monitor's stack confusion all come from the same
weakness: objects stacked or lying on each other in the tray, especially brown on brown.
Each error alone is mostly harmless; the no-movement pre-check absorbs most sub-task false
fails. They become violations when a corrective step acts on them. The replan and the
retry round were both triggered by false verdicts, and both then re-planned without
keeping the prompt's constraint. Of the 8 replan-OFF violations, 4 are this chain: a
final-evaluation misread followed by a retry round that planned an object the prompt
excluded.

\paragraph{Round-1 planning is reliable but not complete.} The other remaining violations
come from the first decomposition: ``after the toy'' was copied into the sub-task text
(2/2 trials of that prompt), ``one of the brown objects'' got two sub-tasks, and ``the
yellow object'' selected all three objects. These are language-understanding failures of
the planner, not perception failures. They concentrate on phrasings that differ from the
planner's examples, which suggests targeted prompt examples or a code-side check of the
sub-task format and count.

\paragraph{Corrective loops need to be conservative.} Turning off in-round replanning
removed a source of violations without losing recoveries. The same holds for retry rounds:
offline, round-2 plans improved from 9/19 to 14/19 correct when the ``judged incomplete''
framing was replaced with a neutral one (``these sub-tasks are done; plan the instruction
against the current scene; nothing left is a valid answer''). Adding a requirements guard
and a count cap gave 16/19. A fresh identify pass did \emph{worse} (6/19): it brought back
invented objects. That retry-round change is not part of the measured system. It is the
obvious next step, and it targets 4 of the 8 remaining violations.

\paragraph{Naming is a hidden interface between the VLM and the VLA.} The VLM calls the
pouch ``case'', ``pencil case'', ``cylinder'' or ``brown tube''. The VLA was trained on one
vocabulary and does not move for unfamiliar names. The planner's exclusion check misses
``case'' for ``except the pouch'' (\#45), and the final-evaluation survey renames objects
between rounds (\#50). A single object list per task, shared by the planner, monitor and
evaluator, would remove a class of errors that never shows up as a perception failure.

\paragraph{Limitations.}
\begin{itemize}
  \item One scene and three objects. 57 trials, mostly one run per (prompt, layout), so
    the per-category numbers are small.
  \item One operator labelled every trial, from live observation.
  \item The replan-OFF set combines first-pass trials with 11 re-runs recorded later in
    the same session.
  \item Configuration settled during the run: trial \#1 used three sub-task retries and
    16-frame clips, and 34 trials ran with replanning enabled but never triggered it.
  \item Two ``complete'' outcomes depend on operator aborts (\#54, \#58).
  \item The baseline gets one 60\,s episode while the system gets several. The comparison
    is between the policy alone and the policy as deployed in the system, not at matched
    compute; the system takes about $2.2\times$ the baseline's wall time.
\end{itemize}

\paragraph{Next steps.}
\begin{enumerate}
  \item Neutral retry-round framing with the requirements guard and count cap; it is
    validated offline (16/19) and targets the retry-round violations.
  \item One object vocabulary per task, shared across the planner, monitor and evaluator.
  \item A post-homing still for the sub-task evaluator, or removing frozen frames from its
    clip, to cut pouch false fails.
  \item Check round-1 plans in code: one object per sub-task in the standard ``put the X on
    the tray'' form, and the number of sub-tasks within any count in the prompt (\#20,
    \#32, \#52).
  \item Guard vocabulary refinement against swapping in another object's name (\#17,
    \#34).
\end{enumerate}
